\subsection{Equivalent Systems and Allowed Row Operations}

Two systems are equivalent when they have the same solution set. Gaussian elimination uses row operations that replace the written equations by equivalent conditions.

\begin{definitionbox}
The allowed operations are: exchanging two rows, multiplying a row by a nonzero scalar, and adding a multiple of one row to another.
\end{definitionbox}

\begin{interpretationbox}[Why row operations are legitimate]
Each operation can be reversed. Therefore it changes the form of the equations without changing which assignments satisfy them.
\end{interpretationbox}

\begin{remarkbox}
The aim is not to make the matrix look different. The aim is to expose the solution set by simplifying the conditions.
\end{remarkbox}
